\documentclass{ximera}
\title{Precise definition of limit}
\begin{document}
\begin{abstract}
The epsilon-delta definition of a limit and the structure of an epsilon-delta proof. Textbook §2.4.
\end{abstract}
\maketitle

Let's be a little more precise by what we mean by a limit.

\begin{definition}[Precise definition of a limit]\label{definition:2-17}
% AUTHOR CHOICE: the definition was blank space in the original (written live); standard textbook statement supplied, change freely
Let $f$ be a function defined on some open interval that contains the number $a$, except possibly at $a$ itself. We say that the \dfn{limit of $f(x)$ as $x$ approaches $a$ is $L$}, and we write
\[
\lim_{x \rightarrow a} f(x)=L,
\]
if for every number $\varepsilon>0$ there is a number $\delta>0$ such that
\[
\text{if } 0<|x-a|<\delta \text{ then } |f(x)-L|<\varepsilon .
\]
\end{definition}

What does this mean?
\begin{image}
\includegraphics[width=0.8\textwidth]{b48a5324-93dc-45ee-b7c2-dc0b5706bda6-023_765_1048_1345_537}
\end{image}

We can use this definition as a way to show whether or not the intuitive solution to a limit problem is accurate or not. In order to prove that our intuition is correct, we use what's called an \emph{epsilon-delta proof}. The following is a general outline of how an epsilon-delta proof works:

\section{Epsilon-delta proof}
% AUTHOR CHOICE: the outline entries were blank lines in the original; wording supplied from the structure of Example 2.18, change freely
\begin{enumerate}
  \item \textbf{Analysis:}
  \begin{enumerate}
    \item Start with $|f(x)-L|<\varepsilon$ and simplify it.
    \item End with something of the form $|x-a|<(\text{an expression in } \varepsilon)$; this tells us what $\delta$ should be.
  \end{enumerate}
  \item \textbf{Proof:}
  \begin{enumerate}
    \item Suppose $\varepsilon>0$.
    \item Let $\delta$ be the expression found in the analysis.
    \item If $0<|x-a|<\delta$, then show that $|f(x)-L|<\varepsilon$.
  \end{enumerate}
  \item Therefore, by the definition of a limit, $\lim_{x \rightarrow a} f(x)=L$.
\end{enumerate}

\begin{example}[Epsilon-delta proof]\label{example:2-18}
% AUTHOR CHOICE: statement to prove chosen during conversion, change freely
We want to prove that $\lim_{x \rightarrow 3}(4x-5)=7$.
\begin{enumerate}
  \item \textbf{Analysis:}
  \begin{enumerate}
    \item We start with $|f(x)-L|<\varepsilon$. In our case, $L=\answer{7}$ and $f(x)=\answer{4x-5}$. Therefore
    \begin{freeResponse}\end{freeResponse}
    % work: |4x-5-7| = |4x-12| = 4|x-3| < epsilon  <=>  |x-3| < epsilon/4
    \item We want to end with something like $|x-3|<\answer{\frac{\varepsilon}{4}}$.
  \end{enumerate}
  \item \textbf{Proof:}
  \begin{enumerate}
    \item Suppose \wordChoice{\choice[correct]{$\varepsilon>0$}\choice{$\delta>0$}\choice{$x=3$}}.
    \item Let $\delta=\answer{\frac{\varepsilon}{4}}$.
    \item If \wordChoice{\choice[correct]{$0<|x-3|<\delta$}\choice{$|f(x)-7|<\varepsilon$}\choice{$0<|x-3|<\varepsilon$}}, then
    \begin{freeResponse}\end{freeResponse}
  \end{enumerate}
  \item Therefore, by the definition of a limit
  \begin{freeResponse}\end{freeResponse}
\end{enumerate}
\end{example}

\begin{exercise}\label{exercise:2-19}
% AUTHOR CHOICE: statement to prove chosen during conversion, change freely
With a partner, prove that $\lim_{x \rightarrow 2}(3x+1)=7$.
\begin{freeResponse}\end{freeResponse}
\end{exercise}

\begin{example}\label{example:2-20}
% AUTHOR CHOICE: statement to prove chosen during conversion, change freely
Let's do a significantly harder example. Prove that $\lim_{x \rightarrow 3} x^{2}=9$.
\begin{enumerate}
  \item \textbf{Analysis:}
  \begin{enumerate}
    \item \begin{freeResponse}\end{freeResponse}
    \item \begin{freeResponse}\end{freeResponse}
  \end{enumerate}
  \item \textbf{Proof:}
  \begin{enumerate}
    \item \begin{freeResponse}\end{freeResponse}
    \item \begin{freeResponse}\end{freeResponse}
    \item \begin{freeResponse}\end{freeResponse}
  \end{enumerate}
  \item \begin{freeResponse}\end{freeResponse}
\end{enumerate}
\end{example}
\end{document}
